\section{Two raw-kernel verifications with endogenous primary laws}
\label{sec:realizations}
Both classes below use the same physical gate/acquisition protocol, but their primary report laws depend on their hidden conditional states. Consequently neither has the exact common finite primary driver of \cite{GTF53}. Let $\nu_0(dy,dt)$ be a specified probability law on compact $K\times[0,1]$, let $|\phi_a|\le1$ with $\int\phi_a\,d\nu_0=0$, and fix $0\le\beta<1$. Put
\[
 c=1-\beta,\quad C=1+\beta,\quad
 0<q\le\tfrac12,\quad 0<\gamma<\tfrac12,\quad
 a_0=\frac{qc\gamma}{C(1-\gamma)}\le\tfrac12.
\]
The functions, bounds and a reference model are known; the actual transition, gate and weak-observation parameters need not be known. All admitted actions obey the same bounds. The audit carrier is chosen compact and convex to contain $y+t e\,e_1$ for $y\in K$, $0\le t\le1$ and $e=\pm1$.

\begin{theorem}[Controlled hidden-state sensing]\label{thm:classical}
Let a hidden sign evolve on every call by a controlled stochastic matrix whose two entries in each row lie in $[q,1-q]$. On a gate call its state-dependent success probabilities lie in $[p_-,p_+]$. On an acquisition call, after that transition, the primary mark has conditional density
\[
 L_a(y,t\mid e)=1+\beta e\phi_a(y,t)
\]
relative to $\nu_0$. An auxiliary weak sign report $v\in\{-1,1\}$ is then produced with each likelihood in $[\gamma,1-\gamma]$. The audit is
\[
 Z=Y+TEe_1,
\]
with the current hidden sign and no extra transition. Then $f(y,t)=y$, $w(y,t)=t^2$, $B_0=0$ satisfy \eqref{eq:ordinary-calibration} with
\begin{equation}\label{eq:raw-calibration}
 a_v=4a_0(1-a_0),\qquad A_v=a_b=1.
\end{equation}
An actually acquired reference law $\nu$ gives conditional envelope constants $\ell=c/C$, $u=C/c$. Thus Theorems~\ref{thm:main} and \ref{thm:herald} apply to every legal history-dependent competitor and one common two-label upper program. If $\nu_0$ has the marked rectangle geometry, Corollary~\ref{cor:regular} gives its matching curve with the same $d,\alpha$.
\end{theorem}
\begin{proof}
Immediately before the primary acquisition report, both sign probabilities are at least $q$, regardless of how much information was obtained or rejected earlier. Since $c\le L_a\le C$, the primary mark's conditional density lies in $[c,C]$ relative to $\nu_0$. After also observing $v$, the numerator of \emph{each} sign posterior is at least $qc\gamma$, while the report likelihood denominator is at most $C(1-\gamma)$. Both posteriors therefore lie in $[a_0,1-a_0]$. In particular $a_0\le1/2$, and if $b=\E[E\mid h,y,t,v]$, then
\[
 \operatorname{Var}(Z\mid h,y,t,v)=(1-b^2)t^2
        \in[4a_0(1-a_0)t^2,t^2],\qquad
 \|\E[Z\mid h,y,t,v]-y\|^2=b^2t^2\le t^2.
\]
This is \eqref{eq:raw-calibration} after the entire report. A fresh stopping coin independent of the audit does not change the conditional identity.

Run the declared gate/acquisition protocol once in a known reference member of the family. Its first productive mark law $\nu$ is an actual acquisition law; conditional densities and their mixtures show $c\nu_0\le\nu\le C\nu_0$. Hence every conditional productive law in every model is between $(c/C)\nu$ and $(C/c)\nu$. This comparison does not supply the true model to the upper program or condition away a tail event. It also transfers rectangle estimates from $\nu_0$ to $\nu$.

The raw hidden transition may have different rows, and gate and primary likelihoods may depend on the sign. Its posterior changes after failed gates and auxiliary reports. In that case both future gate probabilities and primary mark densities change with the observed history, as in Example~\ref{ex:closure}. All preceding estimates remain valid. The program stores its physical phase, not this varying posterior. Normalization denominators on acquisition are bounded below by $c\gamma$ relative to $\nu_0$; thus the report and posterior maps have explicit uniform continuity bounds on the compact pre-report simplex. No contraction or resetting of the hidden process was used in the transfer proof.
\end{proof}

\begin{theorem}[Noncommuting quantum heralding]\label{thm:quantum}
Let the hidden system be a qubit. On each call first apply
\[
 \mathcal D_q(\rho)=(1-2q)\rho+q\tr(\rho)I,
\]
then a controlled unitary $U_a$. Gate effects $G_a,I-G_a$ satisfy $p_-I\le G_a\le p_+I$ and use their positive square roots as Kraus operators. On an acquisition call the primary instrument density is $L_a(y,t)=F_a(y,t)^{1/2}$, where
\[
 F_a(y,t)=I+\beta\phi_a(y,t)n_a\cdot\sigma,
 \qquad \int F_a\,d\nu_0=I.
\]
It is followed by weak effects $E_{a,v}=(I+v\eta m_a\cdot\sigma)/2$, $\eta=1-2\gamma$, and Kraus operators $K_{a,v}=E_{a,v}^{1/2}$. The axes need not commute. The audit is $Z=Y+Te\,e_1$, with $e$ the outcome of a subsequent $\sigma_z$ measurement and no intervening channel. Then the same constants \eqref{eq:raw-calibration} and the same conditional envelopes verify Theorems~\ref{thm:main} and \ref{thm:herald}.
\end{theorem}
\begin{proof}
For a normalized input, the state after channel and unitary satisfies $\rho'\ge qI$. The primary effect obeys $cI\le F_a\le CI$. Its density $\tr(F_a\rho')$ lies in $[c,C]$, and the integrated instrument is trace preserving. After the primary and weak reports the unnormalized state is
\[
 \omega=K_{a,v}L_a\rho' L_aK_{a,v}\ge qc\gamma I.
\]
For normalization,
\[
 \tr\omega=\tr(E_{a,v}L_a\rho' L_a)
       \le(1-\gamma)\tr(F_a\rho')\le C(1-\gamma).
\]
Thus $\omega/\tr\omega\ge a_0I$. Each $\sigma_z$ outcome has probability at least $a_0$, so its mean and variance give precisely the classical audit comparison. The reference acquisition and density envelopes are proved in exactly the same measure-theoretic way, now using $\tr(F_a\rho')\in[c,C]$. Gate success $\tr(G_a\rho')\in[p_-,p_+]$ is generally history dependent. Its state update is not a reset.

Nonparallel $n_a,m_a$ and gate axes yield noncommuting instruments; no shared basis was used. The physical state enters the raw experiment but is not writable or freely readable controller memory. The observer remains a two-live-label classical program. Report-kernel continuity follows from linearity of unnormalized instruments and the positive normalization bounds. All arguments are finite-dimensional. They assert neither unbounded-operator closure nor uniform control at projective or vanishing-$q$ limits.
\end{proof}

\begin{lemma}[Raw report continuity and normalized update stability]\label{lem:raw-continuity}
In either realization, at a productive call the complete primary/weak report law is nonexpansive in total variation with respect to the input state distance, up to the conventional factor $1/2$. For a fixed positive report its normalized predictive update satisfies
\[
 \|F_{a,x,v}(z)-F_{a,x,v}(z')\|_1
 \le \frac{2C(1-\gamma)}{c\gamma}\|z-z'\|_1,
\]
where the norm is total coordinate variation classically and trace norm quantum mechanically. This bound is uniform over the declared acquisition family. It is a one-step stability estimate, not a contraction assertion for arbitrary long products.
\end{lemma}
\begin{proof}
Write the unnormalized report map as $A$. Positivity gives
$\|A(z-z')\|_1\le C(1-\gamma)\|z-z'\|_1$ by the positive/negative decomposition (or its finite-dimensional matrix analogue). Its trace or total mass is at least $c\gamma$ on normalized inputs. Subtract $Az/\tr Az$ and $Az'/\tr Az'$; one term is the unnormalized difference divided by the first denominator, and the other is bounded by the same quantity because the second normalized output has norm one. This proves the displayed factor two. Integrating the absolute report-mass differences uses the trace-preserving complete instrument and gives the total-variation assertion. The preceding transition/channel and unitary are nonexpansive in these norms. No derivative of a posterior at a zero-evidence event is taken.
\end{proof}

\subsection*{Singular marks and the sharp positivity boundary}
Take $\nu_0$ on a Cantor location support with reliability law $\Pp(T\le s)=s^\alpha$. Bounded likelihood factors preserve its actual singular support and rectangle masses; the two raw proofs require no Lebesgue density in ambient coordinates. Alternatively add a marked atom $(y_*,0)$ to the preparation/report instrument, preserving its actual mixture weight in the reference law. Theorems~\ref{thm:main} and \ref{thm:atomic} then retain the revealing transition. The statement is not uniform under arbitrarily changing supports.

The explicit dependence on positivity is $a_0=qc\gamma/[C(1-\gamma)]$, followed by $a_v=4a_0(1-a_0)$ and the constants in \eqref{eq:ordinary-calibration}. Proposition~\ref{prop:q-sharp} shows that this deterioration is necessary for a current-mark calibration over the full admitted family. Uniformity in rare gate success with bounded $p_+/p_-$ must therefore not be confused with uniformity in $q\downarrow0$ or $\gamma\downarrow0$.

\subsection*{A conventional Gaussian audit without a supplied reliability coordinate}
The preceding marked models make the geometry transparent. The next consequence uses the ordinary task of predicting the current hidden sign; its reliability is derived from the raw likelihood, not supplied as an additional random coordinate. Its Gaussian tail mechanism is inherited from the stopped-experiment companion~\cite{GTF52}; the new application pays for history-dependent heralding through the same general theorem.

\begin{theorem}[Endogenously gated Gaussian prediction]\label{thm:gaussian-herald}
Use the two physical ports of Theorem~\ref{thm:herald}. A hidden sign evolves on every call with controlled transition entries in $[q,1-q]$, $0<q\le1/2$. Gate probabilities may depend on its hidden state and lie in $[p_-,p_+]$, with $p_+/p_-\le\Gamma$. A productive call reports
\[
 Y=aE+\sigma G,\qquad G\sim N(0,1),\qquad a,\sigma>0,
\]
with fresh $G$, and the subsequent audit is simply $Z=E$, without an extra transition. The transition and gate parameters are inaccessible. Suppose the known reference family includes a fair post-transition sign and a sign-independent gate. Put $B_0=0$, $\kappa_+=1+1/p_-$ and
\[
 F_G(b)=\exp\!\left[-\frac{2a}{\sigma}\sqrt{2\log b}\right],\qquad b\ge1.
\]
For every $N\ge\kappa_+$, every $m\ge1$ and every $M\ge3$,
\[
 R_\cp^\Theta(N,m)\asymp F_G(N/\kappa_+),\qquad
 R_\aut^\Theta(N,M)\asymp F_G(N/\kappa_+).
\]
Constants depend on $q,a,\sigma,\Gamma$, not on $N,m,M,p_-$. The autonomous upper uses two live states and one terminal prediction. The lower holds for unrestricted history-based stopping and arbitrarily large finite dictionaries. The rare-herald statement is on $N\ge\kappa_+$, not below that feasible comparison range.
\end{theorem}
\begin{proof}
Before the productive Gaussian report the hidden positive probability $r$ lies in $[q,1-q]$. For $y\ge0$, set $t=e^{-2ay/\sigma^2}$ and $c_q=q/(1-q)$, $C_q=1/c_q$. The conditional negative/positive odds are $[(1-r)/r]t\in[c_qt,C_qt]$. Thus for $b_y=\E[E\mid h,Y=y]$,
\[
 \frac{4c_q}{(1+C_q)^2}t\le1-b_y^2\le4C_qt,
 \qquad |b_y-1|^2\le4C_q^2t.
\]
The symmetric calculation applies for $y<0$. Hence the raw task admits the uniform mark
$f(y)=\operatorname{sgn}y$, $w(y)=e^{-2a|y|/\sigma^2}$ by \eqref{eq:ordinary-calibration}. The tie at zero has probability zero. The actual reference acquisition law is the fair mixture
$\nu=\frac12N(a,\sigma^2)+\frac12N(-a,\sigma^2)$. Every full-history productive conditional law lies between $2q\nu$ and $2(1-q)\nu$. This is not a report-row equality: the mixture coefficient $r$ depends on history.

We show uniformly for every $m\ge1$ that $\Qm_m(b)\asymp F_G(b)$. Let $z_b$ be the upper-tail threshold for $|Y|$ with $\nu(|Y|>z_b)=1/(2b)$. Any density $h\le b$ of mass one puts at least half its mass in $|Y|\le z_b$. Discarding the nonnegative codepoint term gives
\[
 \int d_Ch\,d\nu\ge\tfrac12e^{-2az_b/\sigma^2}.
\]
For $b\ge2$, choose $y_b\ge0$ with $\nu(Y\ge y_b)=1/b$, take $h=b\one_{[y_b,\infty)}$ and the single codepoint $+1$. On this set $f=+1$, so the profile is at most $e^{-2ay_b/\sigma^2}$.

For completeness, integration by parts for the standard Gaussian tail gives, for $x>0$,
\[
 \frac{x}{1+x^2}\frac{e^{-x^2/2}}{\sqrt{2\pi}}
 \le\Pp(G\ge x)\le
 \frac1x\frac{e^{-x^2/2}}{\sqrt{2\pi}}.
\]
The upper follows by replacing $1$ in the integral by $u/x$; integration by parts and $u^{-2}\le x^{-2}$ give the lower. These inequalities applied to the two fixed shifted tails imply
\[
 y_b=a+\sigma\sqrt{2\log b}+O\!\left(\frac{\log\log b}{\sqrt{\log b}}\right),\quad
 z_b=a+\sigma\sqrt{2\log b}+O\!\left(\frac{\log\log b}{\sqrt{\log b}}\right)
\]
as $b\to\infty$. Consequently both exponentials are comparable to $F_G(b)$. The remaining bounded $b$ interval has positive continuous quantiles and is handled with constants independent of $m$; a fixed output at the first productive report gives its upper. Finally apply Theorem~\ref{thm:main}. Fixed multiplicative changes of $b$ only change $F_G$ by bounded factors, and $\kappa_+/\kappa_-\le\Gamma$. Two incompatible physical ports and a fixed $+1$ terminal state implement the tail rule. No posterior, hidden transition parameter, or time counter is stored. The mathematical threshold depends only on the known Gaussian and resource certificates; finite-precision evaluation of it requires its own reported numerical accuracy and budget reserve.
\end{proof}
